% =====================================================================
%  BÖLÜM 0 — Ön Bilgiler: Kümeler, Fonksiyonlar ve Sayılabilirlik
%  Olasılık Teorisi ve Matematiksel İstatistik
%  Kaynaklar: Billingsley Ek A; Klenke Böl. 1; Durrett Ek A
% =====================================================================

\chapter{Ön Bilgiler: Kümeler, Fonksiyonlar ve Sayılabilirlik}

\bolumalinti{``Mathematicians always strive to confuse their audiences;
  where there is no confusion, there is no prestige.''}{Gian-Carlo Rota (1932--1999)}

\noindent\textbf{Ön Koşullar:} Lise düzeyinde kümeler ve fonksiyonlar kavramı.
Bir önceki bölüm yoktur; bu bölüm tüm kitabın temel taşıdır.

\noindent\textbf{Bu Bölüm Sonraki Bölümlerde Kullanılır:}
Bölüm~1 (ölçüm teorisi için doğrudan),
Bölüm~2 (olasılık uzayı yapısı),
Bölüm~3 (rasgele değişkenler).

\noindent\textbf{Tahmini Okuma Süresi:} $\approx 4$~saat.

\vspace{0.5\baselineskip}

% ─── KAZANIMLAR ────────────────────────────────────────────────────
\begin{kazanimlar}
Bu bölümün sonunda öğrenci:
\begin{itemize}[leftmargin=1.8em, itemsep=2pt]
  \item Küme cebiri işlemlerini (birleşim, kesişim, fark, tümleyen) hem
        sonlu hem sayılabilir hem de keyfi indis kümeleri üzerinde
        uygulayabilmeli ve De~Morgan yasalarını ispatlayabilmeli,
  \item Ters görüntü (\emph{preimage}) kavramını tanımlayabilmeli ve ters
        görüntünün birleşim, kesişim ve tümleme işlemlerini koruyduğunu
        ispatlayabilmeli,
  \item Sonlu, sayılabilir ve sayılamaz kümeleri ayırt edebilmeli;
        $\Q$'nun sayılabilir, $\R$'nin sayılamaz olduğunu ispatlayabilmeli,
  \item Gerçel sayılarda üst sınır (\emph{supremum}) ve alt sınır
        (\emph{infimum}) kavramlarını tanımlayabilmeli ve tamtümlük
        aksiyomunu ifade edebilmeli,
  \item Genişletilmiş gerçel sayı doğrusu $\Rbar$'yi
        ve üzerinde yapılan aritmetik kuralları açıklayabilmeli,
  \item Sayı dizileri için $\limsup$ ve $\liminf$'i tanımlayabilmeli ve
        hesaplayabilmeli,
  \item Küme dizileri için $\limsup$ ve $\liminf$'i yorumlayabilmeli
        (``sonsuz kez'' ve ``eninde sonunda'' olay dili ile ilişkilendirebilmeli).
\end{itemize}
\end{kazanimlar}

% ─── MOTİVASYON ────────────────────────────────────────────────────
\begin{motivasyon}
Olasılık teorisinin kalbi \emph{ölçüm teorisidir} (measure theory) ve ölçüm
teorisi tamamen küme işlemleri üzerine kuruludur. Bir olayın olasılığını
tanımlamak için, önce ``hangi kümelere olasılık atayabiliriz?'' sorusunu
yanıtlamamız gerekir. Bu sorunun yanıtı $\sigma$-cebiri kavramına götürür;
$\sigma$-cebirleri ise sayılabilir birleşim ve kesişim gibi işlemleri
barındırdıklarından, sayılabilirlik kavramını doğrudan olasılık teorisinin
merkezine taşır.

\medskip

\textbf{Köprü kurma.} Lise matematiğinde sonlu küme işlemlerini (Venn
şemaları, De~Morgan yasaları) öğrendiniz. Bu bölümde o bilgiyi \emph{sonsuz
küme aileleri} üzerine genişleteceğiz. Böylece Bölüm~1'de, sonsuz çok olayı
kapsayan bir $\sigma$-cebirini kolaylıkla tanımlayabilecek ve Kolmogorov'un
aksiyomlarını Bölüm~2'de tam anlamıyla anlayabileceksiniz.

\medskip

Bu bölümde şunları öğreneceğiz:
küme cebiri ile rahat çalışmayı ($\S0.1$),
fonksiyonların ters görüntü özelliklerini ($\S0.2$),
sayılabilirlik teorisini ($\S0.3$),
gerçel sayı sisteminin tamtümlük özelliğini ($\S0.4$)
ve sayı/küme dizilerinde $\limsup$/$\liminf$ kavramlarını ($\S0.5$).
\end{motivasyon}


% =====================================================================
\section{Küme Cebiri}
% =====================================================================

\begin{kesif}[Sonsuz birleşim ne demektir?]
Sonlu iki kümenin birleşimini $A \cup B$ olarak biliyorsunuz. Peki
$A_1 \cup A_2 \cup A_3 \cup \cdots$ ifadesini nasıl yazarsınız?
Düşüncelerinizi aşağıya not edin; ardından Tanım~\ref{def:kumecebiri}'yi
okuyun.
\ogrencialani[2.5cm]
\end{kesif}

\begin{tanimgerekce}
Sonsuz birleşimi $\bigcup_{n=1}^\infty A_n$ olarak yazmak sadece notasyon
kolaylığı değildir. Olasılık teorisinde olayların ``en az birinin gerçekleşmesi''
durumu doğal olarak sonsuz birleşim gerektirir. Örneğin bir madeni paranın
sonsuza dek atılması deneyi düşünüldüğünde, ``en az bir kez tura gelme'' olayı
sonsuz birleşimle ifade edilir.
\end{tanimgerekce}

\begin{tanim}[Temel Küme İşlemleri]\label{def:kumecebiri}
Bir $\Omega$ evren kümesi ve $\Omega$'nın alt kümeleri $A, B, (A_n)_{n \geq 1}$
verilsin. Aşağıdaki işlemleri tanımlıyoruz:
\begin{align*}
  \text{Tümleyen:}\quad
    & A^c \coloneqq \{\omega \in \Omega : \omega \notin A\}, \\[2pt]
  \text{Fark:}\quad
    & A \setminus B \coloneqq A \cap B^c
      = \{\omega \in \Omega : \omega \in A,\ \omega \notin B\}, \\[2pt]
  \text{Simetrik fark:}\quad
    & A \triangle B \coloneqq (A \setminus B) \cup (B \setminus A), \\[2pt]
  \text{Sayılabilir birleşim:}\quad
    & \bigcup_{n=1}^\infty A_n
      \coloneqq \{\omega \in \Omega : \omega \in A_n \text{ bazı } n \in \N \text{ için}\}, \\[2pt]
  \text{Sayılabilir kesişim:}\quad
    & \bigcap_{n=1}^\infty A_n
      \coloneqq \{\omega \in \Omega : \omega \in A_n \text{ her } n \in \N \text{ için}\}.
\end{align*}
İki küme \textbf{ayrık} \emph{(disjoint)} adını alır; eğer $A \cap B = \emptyset$.
Bir $(A_n)$ kümeler ailesi \textbf{ikili ayrık} \emph{(pairwise disjoint)} adını alır
eğer $A_i \cap A_j = \emptyset$, $i \neq j$ için.
\end{tanim}

\begin{kritiksorgu}[Keyfi indis kümeleri]
Yukarıdaki tanım $n \in \N$ için yazıldı. Peki $I$ herhangi bir indis
kümesi olsaydı, yani $\bigcup_{i \in I} A_i$, tanım nasıl değişirdi? $I$
sayılamaz olursa ne olur?
\end{kritiksorgu}

Küme cebiri işlemlerinin en önemli özellikleri De~Morgan yasaları aracılığıyla
ifade edilir. Bu yasalar, ölçüm teorisinde $\sigma$-cebirlerinin tümleme altında
kapalı olduğunu ispat etmede kullanılacak temel araçtır.

\begin{teorem}[De~Morgan Yasaları]\label{thm:demorgan}
Bir $(A_n)_{n \geq 1}$ kümeler ailesi için:
\[
  \left(\bigcup_{n=1}^\infty A_n\right)^c = \bigcap_{n=1}^\infty A_n^c
  \qquad \text{ve} \qquad
  \left(\bigcap_{n=1}^\infty A_n\right)^c = \bigcup_{n=1}^\infty A_n^c.
\]
\end{teorem}

\begin{proof}
\textit{Fikir:} ``$A_1 \cup A_2 \cup \cdots$'nin tümleyeni'' demek ``hiçbir $A_n$'de
bulunmamak'' demektir; bu ise tüm $A_n^c$'lerin kesişimidir.

\medskip

\textbf{Adım 1.} İlk eşitliği kanıtlıyoruz. $\omega \in \left(\bigcup_n A_n\right)^c$
olsun. Bu, $\omega \notin \bigcup_n A_n$ demektir; yani \emph{hiçbir} $n$ için
$\omega \in A_n$ değildir. Dolayısıyla \emph{her} $n$ için $\omega \in A_n^c$,
yani $\omega \in \bigcap_n A_n^c$.

\textbf{Adım 2.} Tersine, $\omega \in \bigcap_n A_n^c$ olsun. Her $n$ için
$\omega \in A_n^c$, yani $\omega \notin A_n$. O hâlde $\omega \notin \bigcup_n A_n$,
yani $\omega \in \left(\bigcup_n A_n\right)^c$.

\medskip

\textbf{Adım 3.} İki içerme birden elde edildiğinden
$\left(\bigcup_n A_n\right)^c = \bigcap_n A_n^c$.
İkinci eşitlik, birincisini $A_n^c$ yerine uygulamakla ya da
$A \mapsto A^c$ simetrisiyle elde edilir.
\end{proof}

\begin{ornek}[De~Morgan yasası --- somut uygulama]\label{ex:demorgan}
$\Omega = \R$ ve $A_n = (0, n)$, $n \geq 1$ olsun.
O hâlde $\bigcup_{n=1}^\infty A_n = (0,\infty)$ ve
\[
  \left(\bigcup_{n=1}^\infty A_n\right)^c = (-\infty,0].
\]
Öte yandan $A_n^c = (-\infty,0] \cup [n,\infty)$, dolayısıyla
\[
  \bigcap_{n=1}^\infty A_n^c = (-\infty,0],
\]
çünkü $[n,\infty)$'ların kesişimi boştur. De~Morgan eşitliği doğrulanmış oldu.
\end{ornek}

Küme dizileri için monoton diziler özellikle önemlidir.

\begin{tanim}[Artan ve Azalan Küme Dizileri]\label{def:monotondizi}
$(A_n)_{n \geq 1}$ bir küme dizisi olsun.
\begin{itemize}[leftmargin=1.8em]
  \item Dizi \textbf{artan} \emph{(increasing)} adını alır; eğer $A_1 \subseteq A_2 \subseteq \cdots$.
        Bu durumda $A_n \nearrow$ gösterimi kullanılır ve
        $\lim_n A_n \coloneqq \bigcup_n A_n$ tanımlanır.
  \item Dizi \textbf{azalan} \emph{(decreasing)} adını alır; eğer $A_1 \supseteq A_2 \supseteq \cdots$.
        Bu durumda $A_n \searrow$ gösterimi kullanılır ve
        $\lim_n A_n \coloneqq \bigcap_n A_n$ tanımlanır.
\end{itemize}
\end{tanim}

\begin{mikroozet}[§0.1 özeti]
Temel küme işlemleri sonsuz indis kümeleri üzerine genişletildi.
De~Morgan yasaları tümleme ile birleşim/kesişim arasındaki dönüşümü sağlar.
Monoton küme dizilerinin limiti birleşim ya da kesişim olarak tanımlanır.
\end{mikroozet}


% =====================================================================
\section{Fonksiyonlar ve Ters Görüntü}
% =====================================================================

\begin{motivasyon}[§0.2]
Bir rasgele değişken, tanımı gereği ölçülebilir bir fonksiyondur.
``Ölçülebilir'' sıfatının ne anlama geldiğini anlamak için,
bir fonksiyonun \emph{ters görüntüsünün} küme işlemlerini nasıl
koruduğunu bilmek şarttır. Bu alt bölüm tam olarak bu teknik temel
için vardır.
\end{motivasyon}

\begin{tanim}[Ters Görüntü (Preimage)]\label{def:tergoriuntu}
$f: X \to Y$ bir fonksiyon ve $B \subseteq Y$ olsun.
$f$'nin $B$ altındaki \textbf{ters görüntüsü} \emph{(preimage, inverse image)}:
\[
  f^{-1}(B) \coloneqq \{x \in X : f(x) \in B\}.
\]
$f$'nin tersi mevcut olmasa bile $f^{-1}(B)$ her zaman tanımlıdır.
\end{tanim}

\begin{hataanalizi}[Ters görüntü ile ters fonksiyon]
$f^{-1}(B)$ gösterimi yalnızca $f$ bijektif olduğunda
ters fonksiyonu ifade eder. Genel durumda $f^{-1}(B)$ bir
\emph{küme} işlemidir ve $f$'nin bijektif olmasını gerektirmez.
Örneğin $f(x) = x^2$ ve $B = \{4\}$ ise $f^{-1}(B) = \{-2, 2\}$.
\end{hataanalizi}

Ters görüntünün olasılık teorisini bu kadar kullanışlı kılan temel özellik
şudur: \emph{ters görüntü tüm küme işlemlerini korur.}

\begin{teorem}[Ters Görüntünün Küme İşlemlerini Koruması]\label{thm:tergoriuntu}
$f: X \to Y$ ve $(B_n)_{n \geq 1} \subseteq \mathcal{P}(Y)$ olsun. O zaman:
\begin{enumerate}[label=(\roman*)]
  \item $f^{-1}\!\left(\bigcup_n B_n\right) = \bigcup_n f^{-1}(B_n)$,
  \item $f^{-1}\!\left(\bigcap_n B_n\right) = \bigcap_n f^{-1}(B_n)$,
  \item $f^{-1}(B^c) = \left(f^{-1}(B)\right)^c$ her $B \subseteq Y$ için.
\end{enumerate}
\end{teorem}

\begin{proof}
\textit{Fikir:} Ters görüntü, ``$f(x) \in$'' koşuluyla çalışır; birleşim ve
kesişim de ``bazı / her'' niceleyicileriyle çalışır. Bu iki yapı birebir örtüşür.

\medskip

\textbf{(i)}\; $x \in f^{-1}\!\left(\bigcup_n B_n\right)$
$\iff f(x) \in \bigcup_n B_n$
$\iff$ bazı $n$ için $f(x) \in B_n$
$\iff$ bazı $n$ için $x \in f^{-1}(B_n)$
$\iff x \in \bigcup_n f^{-1}(B_n)$.

\textbf{(ii)}\; $x \in f^{-1}\!\left(\bigcap_n B_n\right)$
$\iff f(x) \in \bigcap_n B_n$
$\iff$ her $n$ için $f(x) \in B_n$
$\iff$ her $n$ için $x \in f^{-1}(B_n)$
$\iff x \in \bigcap_n f^{-1}(B_n)$.

\textbf{(iii)}\; $x \in f^{-1}(B^c)$
$\iff f(x) \in B^c$
$\iff f(x) \notin B$
$\iff x \notin f^{-1}(B)$
$\iff x \in \left(f^{-1}(B)\right)^c$.
\end{proof}

\begin{notkutu}
Teorem~\ref{thm:tergoriuntu}'nun önemi şuradan gelir: Bölüm~1'de
$\sigma$-cebirlerini tanımladığımızda, bir fonksiyonun ``ölçülebilir'' olması tam
olarak, ters görüntüsünün $\sigma$-cebirini $\sigma$-cebirine taşıması anlamına
gelecektir. Yukarıdaki üç özellik bu taşımanın iyi tanımlı olduğunu garanti eder.
\end{notkutu}

\begin{tanim}[Bijektif Fonksiyon Türleri]\label{def:bijektif}
$f: X \to Y$ bir fonksiyon olsun.
\begin{itemize}[leftmargin=1.8em]
  \item $f$ \textbf{enjektif} \emph{(injective, one-to-one)} adını alır;
        eğer $f(x_1) = f(x_2) \Rightarrow x_1 = x_2$.
  \item $f$ \textbf{sürjektif} \emph{(surjective, onto)} adını alır;
        eğer her $y \in Y$ için bir $x \in X$ mevcutsa $f(x) = y$.
  \item $f$ \textbf{bijektif} \emph{(bijective)} adını alır;
        eğer hem enjektif hem sürjektifse.
        Bijektif fonksiyona aynı zamanda \textbf{bire-bir eşleşme}
        \emph{(one-to-one correspondence)} da denir.
\end{itemize}
\end{tanim}

\begin{kendinleifade}
Teorem~\ref{thm:tergoriuntu}'nun (iii) maddesini kendi cümlenizle ifade edin:
``$f^{-1}(B^c) = (f^{-1}(B))^c$ ifadesi ne anlama geliyor?''
Özellikle $B = Y$ ve $B = \emptyset$ özel durumlarını kontrol edin.
\ogrencialani[2cm]
\end{kendinleifade}

\begin{mikroozet}[§0.2 özeti]
Ters görüntü $f^{-1}(B)$ her fonksiyon için tanımlıdır ve bijektiflik
gerektirmez. Ters görüntü birleşim, kesişim ve tümlemeyi korur. Bu özellik
ölçülebilirlik kavramının temel taşıdır.
\end{mikroozet}


% =====================================================================
\section{Sonlu, Sayılabilir ve Sayılamaz Kümeler}
% =====================================================================

\begin{motivasyon}[§0.3]
``Sayılabilir'' kavramı sezgilerimizi yanıltabilir: $\Z$ sonsuz olmasına
karşın ``sayılabilir'', $\R$ ise sonsuz olup aynı zamanda ``sayılamaz''. Bu
fark, $\sigma$-cebirlerinde neden \emph{sayılabilir} birleşim kullandığımızı
(sayılamaz değil) açıklar. Sayılamaz birleşime izin versek, tüm singletonların
birleşimi $\R$'dir; ama bu birleşime Lebesgue ölçüsü atayamazdık.
\end{motivasyon}

\begin{tanim}[Sayılabilirlik]\label{def:sayilabilir}
Bir $A$ kümesi:
\begin{itemize}[leftmargin=1.8em]
  \item \textbf{sonlu} adını alır; eğer $A = \emptyset$ ya da $\{1,\ldots,n\}$
        ile $A$ arasında bir bijeksiyon varsa (bazı $n \in \N$ için).
  \item \textbf{sayılabilir sonsuz} \emph{(countably infinite)} adını alır;
        eğer $\N$ ile $A$ arasında bir bijeksiyon mevcutsa.
  \item \textbf{sayılabilir} \emph{(countable)} adını alır; eğer sonlu ya da
        sayılabilir sonsuzsa.
  \item \textbf{sayılamaz} \emph{(uncountable)} adını alır; eğer sayılabilir değilse.
\end{itemize}
\end{tanim}

\begin{ornek}[Sayılabilir sonsuz kümeler]\label{ex:sayilabilir}
$\Z$ sayılabilir sonsuzdur: $f: \N \to \Z$,
\[
  f(n) = \begin{cases} n/2 & n \text{ çift}, \\ -(n-1)/2 & n \text{ tek} \end{cases}
\]
bijeksiyonu $0,1,-1,2,-2,3,-3,\ldots$ sırasıyla saydırır.

$\N \times \N$ de sayılabilir sonsuzdur: Cantor'un çapraz saydırma düzeniyle
$(1,1),(1,2),(2,1),(1,3),(2,2),(3,1),\ldots$ sırasıyla sayılır.
\end{ornek}

\begin{teorem}[$\Q$ Sayılabilirdir]\label{thm:Qsayilabilir}
Rasyonel sayılar kümesi $\Q$ sayılabilirdir.
\end{teorem}

\begin{proof}
Her $q \in \Q^+$, $q = p/n$ ($p \in \N$, $n \in \N$, $\gcd(p,n) = 1$) biçiminde
tek türlü yazılır. $\Q^+ \subseteq \N \times \N$ gömülüdür ve $\N \times \N$
sayılabilirdir (Örnek~\ref{ex:sayilabilir}). $\Q = \Q^+ \cup \{0\} \cup (-\Q^+)$
üç sayılabilir kümenin birleşimidir; dolayısıyla sayılabilirdir.
\end{proof}

\begin{teorem}[$\R$ Sayılamazdır --- Cantor'un Köşegen Argümanı]
\label{thm:Rsayilamaz}
Gerçel sayılar kümesi $\R$ sayılamazdır.
Daha güçlü bir ifadeyle, $(0,1)$ aralığı bile sayılamazdır.
\end{teorem}

\begin{proof}
\textit{Fikir:} $(0,1)$'in sayılabilir olduğunu varsayarak çelişki türetiriz;
verilen herhangi bir sayımın ``dışarıda bıraktığı'' bir eleman üretiriz.

\medskip

\textbf{Çelişki varsayımı.} $(0,1)$'in sayılabilir olduğunu, yani
$(0,1) = \{x_1, x_2, x_3, \ldots\}$ olduğunu varsayalım.

\textbf{Adım 1.} Her $x_n$'i ondalık gösterimle yazın:
$x_n = 0.a_{n1}\,a_{n2}\,a_{n3}\cdots$, $a_{nk} \in \{0,1,\ldots,9\}$.
(Sona giden $9$'ların tekrarlanmasını önlemek için: $0.999\cdots = 1$ eşitliğini
göz önünde bulundurarak sonlanan ondalıkları tercih edin.)

\textbf{Adım 2.} Yeni bir $y = 0.b_1 b_2 b_3 \cdots$ sayısı inşa edin:
\[
  b_n \coloneqq
  \begin{cases}
    5, & a_{nn} \neq 5, \\
    6, & a_{nn} = 5.
  \end{cases}
\]

\textbf{Adım 3.} $y \in (0,1)$'dir. Öte yandan $y \neq x_n$ her $n$ için,
çünkü $y$ ile $x_n$'in $n$-inci ondalık basamakları farklıdır ($b_n \neq a_{nn}$).

Bu durum, $y$'nin listede yer almadığı anlamına gelir; $(0,1) = \{x_1, x_2, \ldots\}$
varsayımı çelişkilidir.
\end{proof}

\begin{karsiornek}[Sayılabilir kümenin tümleni sayılamaz olabilir]
$\Q$, $\R$ içinde sayılabilirdir; fakat $\Q^c = \R \setminus \Q$
(irrasyonel sayılar) sayılamazdır. Evreni sayılamaz olan bir kümenin
sayılabilir bir alt kümesinin tümleyeni sayılamaz olmak zorundadır.
\end{karsiornek}

\begin{teorem}[Sayılabilir Ailenin Birleşimi]\label{thm:sayilabilirbir}
Eğer her $n \in \N$ için $A_n$ sayılabilir ise,
$\bigcup_{n=1}^\infty A_n$ da sayılabilirdir.
\end{teorem}

\begin{proof}
Her $A_n$ sayılabilir olduğundan $A_n = \{a_{n1}, a_{n2}, a_{n3}, \ldots\}$
olarak yazılabilir. $\bigcup_n A_n$'ın elemanlarını Cantor'un çapraz tarama
düzeniyle sıralıyoruz:
\[
  a_{11},\; a_{12}, a_{21},\; a_{13}, a_{22}, a_{31},\; \ldots
\]
Bu, $\bigcup_n A_n$'dan $\N$'ye enjektif bir fonksiyon verir (tekrarları atlayın);
dolayısıyla $\bigcup_n A_n$ sayılabilirdir.
\end{proof}

\begin{notkutu}
Teorem~\ref{thm:sayilabilirbir}'in tersi doğru değildir: sayılamaz bir küme
sayılamaz çok sayılabilir parçaya bölünebilir. Örneğin $\R$ sayılamaz,
ama her tekil küme $\{x\}$ sonludur.
\end{notkutu}

% Disiplinlerarası bağlantı — TYMM gereksinimi
\begin{tcolorbox}[maarif@base, colframe=maKopru, colback=maKopru!4,
    colbacktitle=maKopru,
    title={Disiplinlerarası bağlantı --- Bilgisayar bilimi ve hesaplanabilirlik}]
Sayılabilirlik kavramı bilgisayar biliminde merkezi bir rol oynar.
Herhangi bir programlama dilinde yazılabilen tüm programlar kümesi sayılabilirdir
(her program sonlu bir karakter dizisidir). Buna karşın $\R \to \R$
fonksiyonları kümesi sayılamazdır. Dolayısıyla ``hesaplanamaz fonksiyonlar'' vardır;
bu Turing'in durma problemi ile doğrudan bağlantılıdır.
\end{tcolorbox}

\begin{mikroozet}[§0.3 özeti]
Sayılabilirlik, bir kümenin $\N$ ile eşleştirilebilmesi anlamına gelir.
$\Q$ sayılabilir, $\R$ sayılamazdır. Sayılabilir ailenin sayılabilir
birleşimi sayılabilirdir. $\sigma$-cebiri tanımında ``sayılabilir birleşim'' koşulu
bu ayrımı işlevsel kılar.
\end{mikroozet}


% =====================================================================
\section{Gerçel Sayılar ve Genişletilmiş Gerçel Sayı Doğrusu}
% =====================================================================

\begin{motivasyon}[§0.4]
Olasılık teorisinde beklentiler, seriler ve integraller kullanılır. Bunların
bazıları $+\infty$ ya da $-\infty$ değer alabilir. Örneğin bir Poisson rasgele
değişkeninin tüm momentleri sonludur; ama bazı ağır kuyruklu dağılımların birinci
momenti bile sonsuz olabilir. Bu yüzden gerçel sayıları $\pm\infty$ ile
genişletmek zorunludur.
\end{motivasyon}

\begin{tanim}[Supremum ve İnfimum]\label{def:supinf}
$A \subseteq \R$, $A \neq \emptyset$ olsun.
\begin{itemize}[leftmargin=1.8em]
  \item $M \in \R$'ye $A$'nın \textbf{üst sınırı} \emph{(upper bound)} denir;
        eğer $a \leq M$ her $a \in A$ için.
  \item $A$'nın \textbf{üst kenar değeri} \emph{(supremum)}, var olan tüm üst
        sınırların en küçüğüdür; $\sup A$ ile gösterilir.
  \item Benzer biçimde \textbf{alt sınır} \emph{(lower bound)} ve
        \textbf{alt kenar değeri} \emph{(infimum)} $\inf A$ tanımlanır.
\end{itemize}
$A$ yukarıdan sınırlı değilse $\sup A = +\infty$; aşağıdan sınırlı değilse
$\inf A = -\infty$ olarak tanımlanır. Boş küme için $\sup \emptyset = -\infty$
ve $\inf \emptyset = +\infty$ konvansiyonu kullanılır.
\end{tanim}

\begin{aksiyom}[Gerçel Sayıların Tamtümlüğü]
\label{ax:tamtumluk}
Yukarıdan sınırlı her boşalmayan $A \subseteq \R$ kümesinin bir
$\sup A \in \R$ üst kenar değeri vardır.
\end{aksiyom}

\begin{notkutu}
$\Q$ bu aksiyomu sağlamaz: $\{q \in \Q : q^2 < 2\}$ kümesinin
$\Q$ içinde supremumu yoktur ($\sqrt{2} \notin \Q$).
Tamtümlük aksiyomu, $\R$'yi $\Q$'dan ayıran temel özelliktir ve
analizin tüm önemli teoremlerinin (Bolzano-Weierstrass, Heine-Borel, vb.) temelidir.
\end{notkutu}

\begin{tanim}[Genişletilmiş Gerçel Sayı Doğrusu]\label{def:Rbar}
$\Rbar \coloneqq \R \cup \{-\infty, +\infty\}$ kümesi
aşağıdaki düzen ve aritmetik kurallarla donatılır:
\begin{gather*}
  -\infty < x < +\infty \quad \text{her } x \in \R \text{ için},\\[4pt]
  x + \infty = +\infty, \quad
  x - \infty = -\infty \quad (x \in \R), \\[2pt]
  x \cdot (+\infty) = \operatorname{sgn}(x) \cdot \infty \quad (x \neq 0), \\[2pt]
  0 \cdot (\pm\infty) = 0
  \quad \text{\emph{(olasılık teorisinde bu konvansiyon kullanılır!)}}, \\[2pt]
  (+\infty) + (-\infty) \quad \text{tanımsız.}
\end{gather*}
\end{tanim}

\begin{uyari}
$(+\infty) - (+\infty)$ ve $(+\infty) + (-\infty)$ tanımsızdır.
Bu ifadelerle karşılaşıldığında problem ifadesini yeniden düzenlemek gerekir.
Olasılık teorisinde beklenti hesaplarken bu durum ortaya çıkabilir
(Bölüm~4'te ele alınacaktır).
\end{uyari}

\begin{hataanalizi}[$0 \cdot \infty = 0$ konvansiyonunun önemi]
Olasılık teorisinde $0 \cdot \infty = 0$ konvansiyonu \emph{kasıtlı olarak}
seçilmiştir: sıfır olasılıklı bir olayın beklentiye katkısının sıfır olmasını
istiyoruz. Gerçel analizin bazı kitapları bu konvansiyonu kullanmaz; bu nedenle
farklı kaynaklardaki tanımları dikkatle karşılaştırın.
\end{hataanalizi}

\begin{ornek}[Supremum ve İnfimum hesabı]
\begin{enumerate}[label=(\alph*)]
  \item $A = \{1 - 1/n : n \in \N\} = \{0, 1/2, 2/3, 3/4, \ldots\}$
        için $\sup A = 1$ (ama $1 \notin A$) ve $\inf A = 0$ ($0 \in A$).
  \item $B = \{n \in \N\} = \{1,2,3,\ldots\}$ için $\inf B = 1 \in B$ ve
        $\sup B = +\infty$.
  \item $C = \emptyset$ için $\sup C = -\infty$ ve $\inf C = +\infty$
        (konvansiyonel tanım).
\end{enumerate}
\end{ornek}

\begin{mikroozet}[§0.4 özeti]
Tamtümlük aksiyomu, her sınırlı kümenin $\R$'de supremum/infimum
sahibi olduğunu garanti eder.
$\Rbar = \R \cup \{-\infty, +\infty\}$,
özellikle $0 \cdot \infty = 0$ konvansiyonuyla, ölçüm ve olasılık teorisinin
standart çalışma alanıdır.
\end{mikroozet}


% =====================================================================
\section{Üst Limit ve Alt Limit}
% =====================================================================

\begin{motivasyon}[§0.5]
Bir dizi yakınsamıyorsa bile $\limsup$ ve $\liminf$ her zaman tanımlıdır
($\Rbar$'da değerleri mevcuttur). Olasılık teorisinde Borel-Cantelli
lemmalarını (Bölüm~7) ifade etmek için küme dizilerinin $\limsup$'ını,
büyük sayılar kanununu anlamak için ise sayı dizilerinin $\limsup$'ını
sıklıkla kullanacağız.
\end{motivasyon}

\subsection{Sayı Dizilerinde Üst ve Alt Limit}

\begin{tanim}[$\limsup$ ve $\liminf$ --- Sayı Dizileri]\label{def:limsup_sayi}
$(a_n)_{n \geq 1} \subseteq \Rbar$ bir dizi olsun. Tanımlarız:
\[
  \limsup_{n \to \infty} a_n
  \coloneqq \inf_{n \geq 1} \sup_{k \geq n} a_k
  = \lim_{n\to\infty} \sup_{k \geq n} a_k,
\]
\[
  \liminf_{n \to \infty} a_n
  \coloneqq \sup_{n \geq 1} \inf_{k \geq n} a_k
  = \lim_{n\to\infty} \inf_{k \geq n} a_k.
\]
İki limit $\Rbar$'da her zaman mevcuttur; çünkü $s_n \coloneqq \sup_{k \geq n} a_k$
azalan, $t_n \coloneqq \inf_{k \geq n} a_k$ artan bir dizidir.
\end{tanim}

\begin{onerme}[$\limsup$ ve $\liminf$'in Temel Özellikleri]\label{prop:limsup_ozellik}
Her $(a_n) \subseteq \Rbar$ için:
\begin{enumerate}[label=(\roman*)]
  \item $\liminf_n a_n \leq \limsup_n a_n$,
  \item $\lim_n a_n$ mevcutsa (uzun anlamda $\Rbar$'da)
        $\lim_n a_n = \liminf_n a_n = \limsup_n a_n$.
\end{enumerate}
\end{onerme}

\begin{proof}
(i) $n$ sabit alındığında $\inf_{k \geq n} a_k \leq \sup_{k \geq n} a_k$.
$n \to \infty$ alındığında sol tarafın limiti $\liminf$, sağ tarafın limiti
$\limsup$'tır; eşitsizlik korunur.

(ii) $\lim_n a_n = L$ olsun. Her $\varepsilon > 0$ için yeterince büyük
$n$'de $|a_k - L| < \varepsilon$ (her $k \geq n$). Bu,
$\sup_{k \geq n} a_k \leq L + \varepsilon$ ve
$\inf_{k \geq n} a_k \geq L - \varepsilon$ verir; $\varepsilon \to 0$
ile sonuç elde edilir.
\end{proof}

\begin{ornek}[$\limsup$ ve $\liminf$ hesabı]
$a_n = (-1)^n$ için: $\sup_{k \geq n} a_k = 1$ ve $\inf_{k \geq n} a_k = -1$
tüm $n$ için, dolayısıyla
$\limsup_n a_n = 1$ ve $\liminf_n a_n = -1$.
Dizi yakınsamadığından $\lim_n a_n$ mevcut değil, ama $\limsup/\liminf$
her zaman tanımlıdır.
\end{ornek}

\subsection{Küme Dizilerinde Üst ve Alt Limit}

Bu alt bölüm Bölüm~7'deki Borel-Cantelli lemmalarının doğrudan ön koşuludur.

\begin{tanim}[$\limsup$ ve $\liminf$ --- Küme Dizileri]\label{def:limsup_kume}
$(A_n)_{n \geq 1}$ bir küme dizisi olsun. Tanımlarız:
\begin{align*}
  \limsup_{n\to\infty} A_n
  &\coloneqq \bigcap_{n=1}^\infty \bigcup_{k=n}^\infty A_k, \\[4pt]
  \liminf_{n\to\infty} A_n
  &\coloneqq \bigcup_{n=1}^\infty \bigcap_{k=n}^\infty A_k.
\end{align*}
\end{tanim}

\begin{tanimgerekce}
Bu tanımların sezgisel yorumları aşağıdaki önerme ile netleşir.
``Sonsuz kez'' ve ``eninde sonunda'' dili, olasılık teorisindeki olayların
geleceğe dair yorumunu doğrudan yansıtır.
\end{tanimgerekce}

\begin{onerme}[Küme $\limsup$ ve $\liminf$'in Yorumu]\label{prop:limsup_kume}
$(A_n)$ bir küme dizisi olsun. O zaman:
\begin{align*}
  \limsup_{n} A_n
  &= \{\omega : \omega \in A_n \text{ sonsuz sayıda } n \text{ için}\}, \\[4pt]
  \liminf_{n} A_n
  &= \{\omega : \omega \in A_n \text{ eninde sonunda her } n \text{ için}\}.
\end{align*}
Olasılık literatüründe bu kümeler sırasıyla ``s.s.~çok'' (\emph{infinitely often}, i.o.)
ve ``e.s.'' (\emph{eventually}) olay kümeleri olarak adlandırılır.
\end{onerme}

\begin{proof}
$\omega \in \limsup_n A_n = \bigcap_n \bigcup_{k \geq n} A_k$ olsun.
Bu, her $n$ için $\omega \in \bigcup_{k \geq n} A_k$ anlamına gelir; yani
her $n$ için bazı $k \geq n$ ile $\omega \in A_k$. Dolayısıyla sonsuz
sayıda $k$ için $\omega \in A_k$ (yoksa bir $n_0$ bulunurdu ki $k \geq n_0$
için $\omega \notin A_k$, bu $\omega \notin \bigcup_{k \geq n_0} A_k$,
çelişki).

Tersine, $\omega$ sonsuz sayıda $n$'de $A_n$'deyse, her $N$ için $N$'den büyük
bir $k$ ile $\omega \in A_k$ vardır; dolayısıyla $\omega \in \bigcup_{k \geq N} A_k$
her $N$ için, yani $\omega \in \bigcap_N \bigcup_{k \geq N} A_k$.

\medskip

$\liminf$ için: $\omega \in \liminf_n A_n = \bigcup_n \bigcap_{k \geq n} A_k$
iff bazı $n_0$ için her $k \geq n_0$'da $\omega \in A_k$,
yani $\omega$ eninde sonunda her $A_n$'dedir.
\end{proof}

\begin{teorem}[$\liminf \subseteq \limsup$]\label{thm:liminf_subseteq}
Her küme dizisi $(A_n)$ için:
\[
  \liminf_{n} A_n \subseteq \limsup_{n} A_n.
\]
Dizi monoton artanlarsa ya da monoton azalanlarsa, $\liminf = \limsup$.
\end{teorem}

\begin{proof}
``Eninde sonunda her $A_n$'de olmak'' kesinlikle ``sonsuz sayıda $A_n$'de olmak''ı
gerektirir; dolayısıyla $\liminf \subseteq \limsup$.

Dizi artanlarsa ($A_n \nearrow$): $\omega \in \limsup_n A_n$ olsun; sonsuz çok $n$
için $\omega \in A_n$, dolayısıyla bazı $n_0$ için $\omega \in A_{n_0}$, ve
$A_{n_0} \subseteq A_k$ her $k \geq n_0$ için. Yani $\omega \in \liminf_n A_n$.
Azalan dizi benzer argümanla ele alınır.
\end{proof}

\begin{ispatiskele}[Teorem~\ref{thm:liminf_subseteq} --- azalan dizi]
Dizi azalanlarsa ($A_n \searrow$) $\liminf = \limsup$ olduğunu kanıtlayın:

\iskeleadim{1}{$A_n \searrow$ olduğundan, her $n$ için $\bigcap_{k \geq n} A_k = \bosluk{A_n}$'deki en küçük eleman, yani $\bigcap_{k \geq n} A_k = \bigcap_{k=1}^\infty A_k$.}

\iskeleadim{2}{$\liminf_n A_n = \bigcup_n \bigcap_{k \geq n} A_k = \bosluk{\bigcap_{k=1}^\infty A_k}$.}

\iskeleadim{3}{Benzer argümanla $\limsup_n A_n = \bosluk{\bigcap_{k=1}^\infty A_k}$, yani iki küme eşittir.}
\end{ispatiskele}

\begin{ornek}[$\limsup$ ve $\liminf$ küme hesabı]\label{ex:limsupkume}
$\Omega = \R$ ve
\[
  A_n =
  \begin{cases}
    (0,1), & n \text{ tek}, \\
    (0,2), & n \text{ çift}.
  \end{cases}
\]
O zaman:
\[
  \limsup_n A_n = (0,2),
  \qquad
  \liminf_n A_n = (0,1).
\]
Çünkü her $x \in (0,2)$ sonsuz sayıda çift indisli $A_n$'de yer alır;
yalnızca $(0,1)$ kümesi eninde sonunda tüm $A_n$'lerde yer alır.
\end{ornek}

\begin{kendinleifade}
$(A_n)$ ile $B_n = A_n^c$ arasındaki $\limsup$ ve $\liminf$ ilişkisini
De~Morgan yasasını kullanarak türetin. Beklenen sonuç:
$\left(\limsup_n A_n\right)^c = \liminf_n A_n^c$.
\ogrencialani[2.5cm]
\end{kendinleifade}

\begin{mikroozet}[§0.5 özeti]
$\limsup_n A_n$ = sonsuz sık olay; $\liminf_n A_n$ = eninde sonunda
gerçekleşen olay. Her zaman $\liminf \subseteq \limsup$. Monoton dizilerde
ikisi eşittir. Bu dil Bölüm~7'deki Borel-Cantelli lemmalarının tam dilidir.
\end{mikroozet}


% =====================================================================
%  BÖLÜM SONU ÖZETİ
% =====================================================================
\section*{Bölüm Özeti}
\addcontentsline{toc}{section}{Bölüm Özeti}

\noindent
\begin{tcolorbox}[maarif@base, colframe=maOzet, colback=maOzet!4,
    colbacktitle=maOzet, title={Bölüm 0 --- Temel Sonuçlar}]
\begin{itemize}[leftmargin=1.8em, itemsep=3pt]
  \item \textbf{De~Morgan:}
        $\left(\bigcup_n A_n\right)^c = \bigcap_n A_n^c$ ve
        $\left(\bigcap_n A_n\right)^c = \bigcup_n A_n^c$.
  \item \textbf{Ters görüntü:} $f^{-1}$ birleşim, kesişim ve tümlemeyi korur
        (Teorem~\ref{thm:tergoriuntu}).
  \item \textbf{Sayılabilirlik:} $\Q$ sayılabilir; $\R$ sayılamaz (Cantor köşegen argümanı).
        Sayılabilir ailenin birleşimi sayılabilirdir.
  \item \textbf{Tamtümlük:} Sınırlı her boşalmayan $A \subseteq \R$ kümesinin
        $\sup A$ ve $\inf A$ değerleri $\R$'de mevcuttur.
  \item \textbf{$\Rbar$:} $0 \cdot (\pm\infty) = 0$ konvansiyonuyla kullanılır;
        $(+\infty)+(-\infty)$ tanımsız.
  \item \textbf{$\limsup / \liminf$ (kümeler):}
        $\limsup_n A_n$ = s.s.~çok; $\liminf_n A_n$ = eninde~sonunda.
        $\left(\limsup_n A_n\right)^c = \liminf_n A_n^c$.
\end{itemize}

\medskip\noindent\textbf{Notasyon özeti:}
\begin{itemize}[leftmargin=1.8em, itemsep=2pt, label={--}]
  \item $A^c$: tümleyen;\quad $A \triangle B$: simetrik fark;\quad $\powerset{\Omega}$: kuvvet kümesi
  \item $f^{-1}(B)$: ters görüntü
  \item $\sup A$, $\inf A$: üst/alt kenar değer;\quad
        $\Rbar = \R \cup \{-\infty, +\infty\}$
  \item $\limsup_n A_n = \bigcap_n \bigcup_{k \geq n} A_k$;\quad
        $\liminf_n A_n = \bigcup_n \bigcap_{k \geq n} A_k$
\end{itemize}
\end{tcolorbox}

% ─── ÖZ-YANSITMA ──────────────────────────────────────────────────
\begin{ozyansima}
Bu bölümü tamamladıktan sonra kendinize şu soruları sorun:
\begin{itemize}[leftmargin=1.8em, itemsep=3pt]
  \item De~Morgan yasasını kaynağa bakmadan ispatlayabilir misiniz?
  \item Ters görüntünün neden tümlemeyi koruduğunu, bir arkadaşınıza sözlü
        açıklayabilir misiniz?
  \item Cantor'un köşegen argümanında hangi adım çelişkiyi üretir?
  \item $\limsup_n A_n$ ile $\liminf_n A_n$ arasındaki farkı tek cümleyle
        ifade edebilir misiniz?
\end{itemize}
\end{ozyansima}

% ─── KÖPRÜ NOTU ───────────────────────────────────────────────────
\begin{kopru}
\textbf{Sonraki bölüm (Bölüm~1: Ölçüm Teorisinin Temelleri).}
Küme cebiri bilgisi doğrudan $\sigma$-cebiri tanımına gidecek;
sayılabilirlik orada ``sayılabilir birleşim altında kapalı'' koşulunda karşınıza
çıkacak. Ters görüntünün küme işlemlerini koruması ise ölçülebilir fonksiyon
tanımının kalbidir.

\medskip

\textbf{İleri okuma.}
Billingsley \cite{billingsley1995} Ek~A;
Klenke \cite{klenke2014} Bölüm~1.1;
Durrett \cite{durrett2019} Ek~A.1.
\end{kopru}


% =====================================================================
%  ALIŞTIRMALAR
% =====================================================================

\cozumlualistirmalar

% ─── Çözümlü Alıştırma 0.1 ───────────────────────────────────────
\begin{alistirma}[Al.0.1]{A}{De~Morgan yasası --- doğrudan uygulama}
$\Omega = \{1,2,3,4,5,6\}$ (zar atma örnek uzayı) ve $A_n = \{k \in \Omega : k \leq n+1\}$,
$n = 1,2,3$ olsun. Hesaplayın:
(a)~$\bigcup_{n=1}^{3} A_n$,\quad
(b)~$\left(\bigcup_{n=1}^{3} A_n\right)^c$,\quad
(c)~$\bigcap_{n=1}^{3} A_n^c$.
İki sonucun eşit olduğunu doğrulayın.
\end{alistirma}
\alcozum{%
(a) $A_1 = \{1,2\}$, $A_2 = \{1,2,3\}$, $A_3 = \{1,2,3,4\}$.
$\bigcup_{n=1}^{3} A_n = \{1,2,3,4\}$.

(b) $\left(\bigcup_n A_n\right)^c = \Omega \setminus \{1,2,3,4\} = \{5,6\}$.

(c) $A_1^c = \{3,4,5,6\}$, $A_2^c = \{4,5,6\}$, $A_3^c = \{5,6\}$.
$\bigcap_n A_n^c = \{5,6\}$.

İki küme eşittir; De~Morgan yasası doğrulanmış oldu. $\blacksquare$%
}

% ─── Çözümlü Alıştırma 0.2 ───────────────────────────────────────
\begin{alistirma}[Al.0.2]{B}{Ters görüntü --- aralıklar}
$f: \R \to \R$, $f(x) = x^2$ olsun.
(a)~$f^{-1}([0,4])$'ü hesaplayın.
(b)~$f^{-1}((-1,0))$'ı hesaplayın ve yorumlayın.
(c)~$B_1 = [1,4]$ ve $B_2 = [4,9]$ için $f^{-1}(B_1 \cap B_2)$'yi
    hem doğrudan hem de $f^{-1}(B_1) \cap f^{-1}(B_2)$ olarak hesaplayın;
    eşit olduklarını doğrulayın.
\end{alistirma}
\alcozum{%
(a) $f^{-1}([0,4]) = \{x : x^2 \in [0,4]\} = [-2,2]$.

(b) $f^{-1}((-1,0)) = \{x : x^2 \in (-1,0)\} = \emptyset$.
Yorum: $f(x) = x^2 \geq 0$ olduğundan görüntü $[0,\infty)$;
negatif değerlerin ters görüntüsü boştur.

(c) $B_1 \cap B_2 = \{4\}$, dolayısıyla $f^{-1}(\{4\}) = \{-2,2\}$.

$f^{-1}(B_1) = [-2,-1] \cup [1,2]$ ve $f^{-1}(B_2) = [-3,-2] \cup [2,3]$.
$f^{-1}(B_1) \cap f^{-1}(B_2) = \{-2,2\}$.

İki sonuç eşittir: Teorem~\ref{thm:tergoriuntu}(ii) doğrulandı. $\blacksquare$%
}

% ─── Çözümlü Alıştırma 0.3 ───────────────────────────────────────
\begin{alistirma}[Al.0.3]{B}{Küme $\limsup$ ve $\liminf$ --- azalan dizi}
$\Omega = [0,1]$ ve $A_n = [0, 1/n]$ olsun.
(a)~$\limsup_n A_n$'i hesaplayın.
(b)~$\liminf_n A_n$'i hesaplayın.
(c)~Dizi monoton mu? Cevabınızı Teorem~\ref{thm:liminf_subseteq} ile
    ilişkilendirin.
\end{alistirma}
\alcozum{%
(a) Dizi azalan: $A_{n+1} \subset A_n$.
$\bigcup_{k \geq n} A_k = A_n = [0,1/n]$ (en büyük eleman $A_n$).
$\limsup_n A_n = \bigcap_{n \geq 1} [0,1/n] = \{0\}$.

(b) $\bigcap_{k \geq n} A_k = \bigcap_{k \geq n} [0,1/k] = \{0\}$ her $n$ için.
$\liminf_n A_n = \bigcup_{n \geq 1} \{0\} = \{0\}$.

(c) Dizi azalan; Teorem~\ref{thm:liminf_subseteq}'ye göre
$\limsup = \liminf = \{0\}$, yani $\lim_n A_n = \{0\}$.
Sezgisel: aralıklar gittikçe küçülür, tek ortak nokta $\{0\}$ kalır. $\blacksquare$%
}

% ─── Çözümlü Alıştırma 0.4 ───────────────────────────────────────
\begin{alistirma}[Al.0.4]{A}{$\Z$'nin sayılabilirliği}
$\Z$'nin sayılabilir sonsuz olduğunu, açık bir bijeksiyon
$f: \N \to \Z$ yazarak kanıtlayın. ($\N = \{1,2,3,\ldots\}$.)
\end{alistirma}
\alcozum{%
$f(n) = \begin{cases} n/2 & n \text{ çift}, \\ -(n-1)/2 & n \text{ tek}\end{cases}$
tanımlayın. Değerler: $f(1)=0$, $f(2)=1$, $f(3)=-1$, $f(4)=2$, $f(5)=-2$, \ldots

\emph{Enjektif:} $f(m) = f(n)$ olsun. $m,n$ ikisi de çift ise $m/2 = n/2$, yani $m=n$.
İkisi de tek ise $(m-1)/2 = (n-1)/2$, yani $m=n$. Biri çift biri tek olamaz
(biri pozitif, diğeri sıfır ya da negatif verir; çelişki).

\emph{Sürjektif:} $z > 0$ ise $n = 2z$ çift, $f(2z) = z$; $z \leq 0$ ise
$n = 1-2z$ tek, $f(1-2z) = z$.

$f$ bijektiftir; dolayısıyla $\Z$ sayılabilir sonsuzdur. $\blacksquare$%
}

\alistirmalar

% ─── A Seviyesi ───────────────────────────────────────────────────
\begin{alistirma}[Al.0.5]{A}{Küme özdeşlikleri}
Aşağıdakilerin doğru mu yanlış mı olduğunu belirleyin. Doğruysa ispatlayın;
yanlışsa karşı örnek verin.
\begin{enumerate}[label=(\alph*)]
  \item $A \setminus (B \cup C) = (A \setminus B) \cap (A \setminus C)$
  \item $A \triangle B = B \triangle A$
  \item $A \cap (B \triangle C) = (A \cap B) \triangle (A \cap C)$
  \item $(A \triangle B) \triangle C = A \triangle (B \triangle C)$
\end{enumerate}
\end{alistirma}
\alcozum{%
(a) Doğru. $A \setminus (B \cup C) = A \cap (B \cup C)^c = A \cap (B^c \cap C^c)
= (A \cap B^c) \cap (A \cap C^c) = (A \setminus B) \cap (A \setminus C)$.

(b) Doğru. Simetrik fark tanımı gereği simetriktir:
$(A \setminus B) \cup (B \setminus A) = (B \setminus A) \cup (A \setminus B)$.

(c) Doğru. Her $\omega$ için dört durum kontrol edilir;
$A \cap (B \triangle C)$ ve $(A \cap B) \triangle (A \cap C)$
her iki ifade ``$A$'da, ama $B$ ve $C$'den tam olarak birinde'' olan
elemanlar kümesidir.

(d) Doğru (simetrik fark birleştirme yasasına uyar). İspat: her $\omega$
için $A \triangle B \triangle C$ tam olarak tek sayıda kümede bulunan
elemanları içerir. $\blacksquare$%
}

\begin{alistirma}[Al.0.6]{A}{Supremum ve İnfimum}
Aşağıdaki kümelerin $\sup$ ve $\inf$ değerlerini bulun; her değerin kümede
bulunup bulunmadığını belirtin.
\begin{enumerate}[label=(\alph*)]
  \item $A = \left\{\dfrac{(-1)^n}{n} : n \in \N\right\}$
  \item $B = \{x \in \Q : x^2 < 3\}$
  \item $C = \left\{1 - \dfrac{1}{2^n} : n \in \N_0\right\}$
        ($\N_0 = \{0,1,2,\ldots\}$)
\end{enumerate}
\end{alistirma}
\alcozum{%
(a) $A = \{-1, 1/2, -1/3, 1/4, \ldots\}$. $\sup A = 1/2 \in A$;
$\inf A = -1 \in A$.

(b) $B = (-\sqrt{3}, \sqrt{3}) \cap \Q$. $\sup B = \sqrt{3} \notin B$
(irrasyonel); $\inf B = -\sqrt{3} \notin B$.

(c) $C = \{0, 1/2, 3/4, 7/8, \ldots\}$. $\sup C = 1 \notin C$
(limit değer, küme içermez); $\inf C = 0 \in C$. $\blacksquare$%
}

\begin{alistirma}[Al.0.7]{A}{$\limsup$ ve $\liminf$ --- salınan dizi}
$(a_n)$ dizisi $a_n = \sin(n\pi/2)$, $n \geq 1$ olsun.
(a)~İlk sekiz terimi yazın.
(b)~$\limsup_n a_n$ ve $\liminf_n a_n$'i hesaplayın.
\end{alistirma}
\alcozum{%
(a) $1,\,0,\,-1,\,0,\,1,\,0,\,-1,\,0$.

(b) Dizi $\{1,0,-1,0,\ldots\}$ döngüsündedir. $\sup_{k \geq n} a_k = 1$ her
$n$ için. $\limsup_n a_n = 1$. Benzer biçimde $\liminf_n a_n = -1$. $\blacksquare$%
}

\begin{alistirma}[Al.0.8]{A}{Gösterge fonksiyonu ve simetrik fark}
$\gostf{A}(\omega) = 1$ eğer $\omega \in A$, aksi hâlde $0$ olsun.
Gösterin ki:
\[
  \gostf{A \triangle B} = \gostf{A} + \gostf{B} - 2\,\gostf{A \cap B}.
\]
\end{alistirma}
\alcozum{%
Dört durum:
\begin{itemize}
  \item $\omega \in A \cap B$: sol = 0, sağ = $1+1-2 = 0$. \checkmark
  \item $\omega \in A \setminus B$: sol = 1, sağ = $1+0-0 = 1$. \checkmark
  \item $\omega \in B \setminus A$: sol = 1, sağ = $0+1-0 = 1$. \checkmark
  \item $\omega \notin A \cup B$: sol = 0, sağ = $0+0-0 = 0$. \checkmark
\end{itemize}
$\blacksquare$%
}

% ─── B Seviyesi ───────────────────────────────────────────────────
\begin{alistirma}[Al.0.9]{B}{Sayılabilir birleşim --- uygulama}
(a)~$\N \times \N$'nin sayılabilir olduğunu ispatlayın.
(b)~$\Q$'nun sayılabilirliğini Teorem~\ref{thm:sayilabilirbir}'i kullanarak
    yeniden ispatlayın.
(c)~Herhangi bir $A \subseteq B$ için $A$ sayılabilir olduğunda $B$'nin
    sayılabilmesinin \emph{yeterli} koşulunu tartışın. ($B$ sayılabilir
    olmak zorunda mı?)
\end{alistirma}
\alcozum{%
(a) $f: \N \times \N \to \N$, $f(m,n) = 2^{m-1}(2n-1)$ bijeksiyondur
(her pozitif tam sayı tek türlü $2^a \cdot b$ biçiminde yazılır, $b$ tek).
Dolayısıyla $\N \times \N \cong \N$, sayılabilirdir.

(b) $\Q^+ \subseteq \N \times \N$ gömülebilir ($p/q \mapsto (p,q)$ enjektif).
$\Q = \Q^+ \cup \{0\} \cup (-\Q^+)$, Teorem~\ref{thm:sayilabilirbir} ile sayılabilir.

(c) $A \subseteq B$ ise $A$'nın sayılabilirliği $B$'yi zorunlu olarak
sayılabilir kılmaz. Örneğin $A = \{0\}$ ve $B = \R$: $A$ sonlu,
$B$ sayılamazdır. Ama $B \subseteq A$ ile $A$ sayılabilir ise $B$
sayılabilirdir (sayılabilir bir kümenin her alt kümesi sayılabilir). $\blacksquare$%
}

\begin{alistirma}[Al.0.10]{B}{Küme $\limsup$ ve $\liminf$ --- iki dizi}
$\Omega = \R$ olsun. Her dizi için $\limsup$ ve $\liminf$'i hesaplayın.
\begin{enumerate}[label=(\alph*)]
  \item $A_n = \left(-\dfrac{1}{n},\, 1 + \dfrac{1}{n}\right)$
  \item $A_n = (n, n+1)$
\end{enumerate}
\end{alistirma}
\alcozum{%
(a) Dizi azalan: $A_{n+1} \subset A_n$. $\bigcup_{k \geq n} A_k = A_n$ ve
$\bigcap_{k \geq n} A_k \to [0,1]$.
$\limsup_n A_n = \liminf_n A_n = [0,1]$.

(b) $\bigcup_{k \geq n} A_k = (n,\infty)$ ve $\bigcap_{k \geq n} A_k = \emptyset$.
$\limsup_n A_n = \bigcap_{n \geq 1}(n,\infty) = \emptyset$.
$\liminf_n A_n = \bigcup_{n \geq 1} \emptyset = \emptyset$.
Yorum: $(n,n+1)$ sağa kayar; hiçbir sabit nokta bu dizinin sonsuz
üyesinde değildir. $\blacksquare$%
}

\begin{alistirma}[Al.0.11]{B}{Ters görüntü --- mutlak değer fonksiyonu}
$f: \R \to \R$, $f(x) = |x|$ olsun.
\begin{enumerate}[label=(\alph*)]
  \item $f^{-1}([1,2])$
  \item $f^{-1}(\{0\})$
  \item $f^{-1}(\R \setminus [0,1])$ hem doğrudan hem de
        $\left(f^{-1}([0,1])\right)^c$ olarak hesaplayın; eşitliği
        Teorem~\ref{thm:tergoriuntu}(iii) ile ilişkilendirin.
\end{enumerate}
\end{alistirma}
\alcozum{%
(a) $f^{-1}([1,2]) = \{x : |x| \in [1,2]\} = [-2,-1] \cup [1,2]$.

(b) $f^{-1}(\{0\}) = \{0\}$.

(c) Doğrudan: $\R \setminus [0,1] = (-\infty,0) \cup (1,\infty)$;
$f^{-1} = (-\infty,-1) \cup (1,\infty)$ (çünkü $|x| > 1 \iff x > 1$ ya da $x < -1$).

$(f^{-1}([0,1]))^c = [-1,1]^c = (-\infty,-1) \cup (1,\infty)$.

İki sonuç eşittir; Teorem~\ref{thm:tergoriuntu}(iii) doğrulandı. $\blacksquare$%
}

\begin{alistirma}[Al.0.12]{B}{$\limsup$ tümleme özdeşliği --- ispat}
$(A_n)$ bir küme dizisi için $\left(\limsup_n A_n\right)^c = \liminf_n A_n^c$
olduğunu ispatlayın.
\end{alistirma}
\alcozum{%
De~Morgan yasasını art arda uygulayın:
\begin{align*}
  \left(\limsup_n A_n\right)^c
  &= \left(\bigcap_{n=1}^\infty \bigcup_{k \geq n} A_k\right)^c \\
  &= \bigcup_{n=1}^\infty \left(\bigcup_{k \geq n} A_k\right)^c \\
  &= \bigcup_{n=1}^\infty \bigcap_{k \geq n} A_k^c \\
  &= \liminf_n A_n^c. \qquad \blacksquare
\end{align*}
Sezgisel yorum: ``$\omega$ sonsuz sık $A_n$'de \emph{değil}'' $\iff$
``$\omega$ eninde sonunda her $A_n^c$'de.'' $\blacksquare$%
}

% ─── C Seviyesi ───────────────────────────────────────────────────
\begin{alistirma}[Al.0.13]{C}{Cantor kümesinin sayılamazlığı}
\textbf{Cantor kümesi} $\mathcal{C}$, $[0,1]$'den ardı ardına orta üçte bir
açık aralıklar çıkarılarak elde edilir: önce $(1/3,2/3)$, ardından
$(1/9,2/9)$ ve $(7/9,8/9)$, vb.

(a)~$\mathcal{C}$'nin Lebesgue ölçüsünün $0$ olduğunu gösterin
(çıkarılan açık aralıkların uzunluk toplamı $1$'dir).

(b)~$\mathcal{C}$, üçlü taban gösteriminde yalnızca $0$ ve $2$ basamağı
içeren sayıların kümesidir. Bu gözlemi kullanarak $\mathcal{C}$'nin
sayılamaz olduğunu kanıtlayın.
\end{alistirma}
\alcozum{%
\begin{ipucu}{1}
$n$-inci adımda $2^{n-1}$ adet, her biri $1/3^n$ uzunluklu aralık çıkarılır.
\end{ipucu}

(a) Toplam çıkarılan uzunluk:
$\sum_{n=1}^\infty 2^{n-1}/3^n = \tfrac{1}{2}\sum_{n=1}^\infty (2/3)^n
= \tfrac{1}{2} \cdot \tfrac{2/3}{1/3} = 1$.

(b) Her $x \in \mathcal{C}$, $x = \sum_{n=1}^\infty d_n/3^n$ ($d_n \in \{0,2\}$)
biçiminde yazılır. $g: \mathcal{C} \to [0,1]$, $g(x) = \sum_{n=1}^\infty (d_n/2)/2^n$
sürjektiftir (her ikili gösterim elde edilir). $[0,1]$ sayılamaz; sürjektif
görüntünün alt kümesi sayılamaz olamaz; dolayısıyla $\mathcal{C}$ sayılamazdır.
$\blacksquare$%
}

\begin{alistirma}[Al.0.14]{C}{Alt dizi ve $\limsup$/$\liminf$ ilişkisi}
$(A_n)$ bir küme dizisi ve $(A_{n_k})_{k \geq 1}$, $(A_n)$'nin bir alt dizisi
($n_1 < n_2 < \cdots$) olsun. Kanıtlayın ki:
\[
  \liminf_n A_n \subseteq \liminf_k A_{n_k}
  \subseteq \limsup_k A_{n_k} \subseteq \limsup_n A_n.
\]
\end{alistirma}
\alcozum{%
\begin{ipucu}{1}
$\limsup_k A_{n_k} \subseteq \limsup_n A_n$ için: $\omega$'nın sonsuz çok $A_{n_k}$'da
olması, sonsuz çok $A_n$'da olması anlamına gelir (alt dizi endeksleri de $\N$'dedir).
\end{ipucu}

\textbf{$\limsup_k A_{n_k} \subseteq \limsup_n A_n$:}
$\omega$ sonsuz çok $k$ için $A_{n_k}$'da $\Rightarrow$ sonsuz çok $n \in \{n_1,n_2,\ldots\}$
için $\omega \in A_n$ $\Rightarrow$ $\omega \in \limsup_n A_n$.

\textbf{$\liminf_n A_n \subseteq \liminf_k A_{n_k}$:}
$\omega \in \liminf_n A_n$ $\Rightarrow$ bazı $N$ için $k \geq N \Rightarrow \omega \in A_k$.
$n_k \geq k$ olduğundan, yeterince büyük $k$ için $n_k \geq N$ ve $\omega \in A_{n_k}$;
yani $\omega \in \liminf_k A_{n_k}$.

Ortadaki $\liminf \subseteq \limsup$ Teorem~\ref{thm:liminf_subseteq}'den gelir. $\blacksquare$%
}

\begin{alistirma}[Al.0.15]{C}{Simetrik fark ve ölçü}
$(X, \mathcal{A}, \mu)$ bir ölçüm uzayı, $A, B \in \mathcal{A}$ ve $\mu$ sonlu
ölçü olsun. Gösterin ki $\mu(A \triangle B) = \mu(A) + \mu(B) - 2\mu(A \cap B)$.
Bu ifade neden bir ``küme metriği'' için temel oluşturur?
\end{alistirma}
\alcozum{%
$A \triangle B = (A \setminus B) \cup (B \setminus A)$ ve bu iki küme ayrık.
\begin{align*}
  \mu(A \triangle B)
  &= \mu(A \setminus B) + \mu(B \setminus A) \\
  &= [\mu(A) - \mu(A \cap B)] + [\mu(B) - \mu(A \cap B)] \\
  &= \mu(A) + \mu(B) - 2\mu(A \cap B).
\end{align*}
Metrik: $d(A,B) = \mu(A \triangle B)$ tanımlanırsa, $d(A,A)=0$, $d(A,B)=d(B,A)$
ve üçgen eşitsizliği $A \triangle C \subseteq (A \triangle B) \cup (B \triangle C)$
özdeşliğinden gelir. Bu ``ölçü teorik metrik'', $L^1$ normunun gösterge
fonksiyonları üzerindeki karşılığıdır. $\blacksquare$%
}
